\subsection{Cardinal Comparison and the Continuum}

Throughout this subsection, assume Principle \ref{princ:axiom-of-choice-analysis} and use the initial-ordinal cardinality $\abs{X}$ from Definition \ref{def:cardinals-as-initial-ordinals}.

\begin{definition} \label{def:cardinal-comparison-and-arithmetic} \hlt{Cardinal Comparison and Arithmetic}\\
For sets $A,B$, define the \hlt{function set} by $A^B\defeq\{f\mid f:B\to A\}$. Let $\kappa,\lambda$ be cardinals.
\begin{enumerate}[label=(\roman*)]
\setlength{\itemsep}{0pt}
\item \hlt{Comparison}: the relations $\kappa\leq\lambda$ and $\kappa<\lambda$ are the corresponding ordinal comparisons.
\item The \hlt{cardinal sum} is $\kappa+\lambda\defeq\abs{(\{0\}\times\kappa)\cup(\{1\}\times\lambda)}$.
\item The \hlt{cardinal product} is $\kappa\lambda\defeq\abs{\kappa\times\lambda}$.
\item The \hlt{cardinal power} is $\kappa^\lambda\defeq\abs{\{f\mid f:\lambda\to\kappa\}}$.
\end{enumerate}
These are cardinal operations. Cardinal arguments distinguish them from the ordinal operations of Definition \ref{def:ordinal-arithmetic}.
\end{definition}

\begin{proposition} \label{prop:cardinal-operations-on-representatives} \hlt{Cardinal Comparison and Operations on Representatives}\\
Let $A,B$ be sets, and set $\kappa\defeq\abs{A}$ and $\lambda\defeq\abs{B}$.
\begin{enumerate}[label=(\roman*)]
\setlength{\itemsep}{0pt}
\item $\kappa\leq\lambda$ iff an injection $A\to B$ exists.
\item $\kappa<\lambda$ iff an injection $A\to B$ exists and $A\not\sim_{\#}B$.
\item $\kappa+\lambda=\abs{(\{0\}\times A)\cup(\{1\}\times B)}$. If $A\cap B=\emptyset$, then $\kappa+\lambda=\abs{A\cup B}$.
\item $\kappa\lambda=\abs{A\times B}$.
\item $\kappa^\lambda=\abs{A^B}$.
\end{enumerate}
\begin{proof}
Choose bijections $u:A\to\kappa$ and $v:B\to\lambda$ by Proposition \ref{prop:cardinality-canonical-and-bijection-invariant}(i).
\begin{enumerate}[label=(\roman*)]
\setlength{\itemsep}{0pt}
\item[$(i)$] ($\Rightarrow$) If $\kappa\leq\lambda$, then an injection $A\to B$ exists by Proposition \ref{prop:ordinal-cardinality-comparison}.
\item[$(i)$] ($\Leftarrow$) If an injection $A\to B$ exists, then $\kappa\leq\lambda$ by Proposition \ref{prop:ordinal-cardinality-comparison}.
\item[$(ii)$] ($\Rightarrow$) Suppose $\kappa<\lambda$. Then $\kappa\leq\lambda$ and $\kappa\neq\lambda$. An injection $A\to B$ exists by (i), and $A\not\sim_{\#}B$ by Proposition \ref{prop:cardinality-canonical-and-bijection-invariant}(ii).
\item[$(ii)$] ($\Leftarrow$) Suppose an injection $A\to B$ exists and $A\not\sim_{\#}B$. Then $\kappa\leq\lambda$ by (i) and $\kappa\neq\lambda$ by Proposition \ref{prop:cardinality-canonical-and-bijection-invariant}(ii), so $\kappa<\lambda$.
\item[$(iii)$] Set $U\defeq(\{0\}\times A)\cup(\{1\}\times B)$ and $V\defeq(\{0\}\times\kappa)\cup(\{1\}\times\lambda)$. Define $F:U\to V$ by $F(0,a)\defeq(0,u(a))$ and $F(1,b)\defeq(1,v(b))$.\\
The inverse sends $(0,\xi)$ to $(0,u^{-1}(\xi))$ and $(1,\eta)$ to $(1,v^{-1}(\eta))$, so $F$ is bijective by Theorem \ref{thm:bijective-iff-invertible-function}.\\
Hence $\abs{U}=\abs{V}$ by Proposition \ref{prop:cardinality-canonical-and-bijection-invariant}(ii), proving the first equality.\\
If $A\cap B=\emptyset$, define $G:U\to A\cup B$ by $G(0,a)\defeq a$ and $G(1,b)\defeq b$.\\
Its inverse sends $x\in A$ to $(0,x)$ and $x\in B$ to $(1,x)$, so $G$ is bijective. Thus $\kappa+\lambda=\abs{A\cup B}$ by Proposition \ref{prop:cardinality-canonical-and-bijection-invariant}(ii).
\item[$(iv)$] Define $P:A\times B\to\kappa\times\lambda$ by $P(a,b)\defeq(u(a),v(b))$. Its inverse is $(\xi,\eta)\mapsto(u^{-1}(\xi),v^{-1}(\eta))$, so $P$ is bijective by Theorem \ref{thm:bijective-iff-invertible-function}.\\
Therefore $\kappa\lambda=\abs{A\times B}$ by Proposition \ref{prop:cardinality-canonical-and-bijection-invariant}(ii).
\item[$(v)$] Define $T:A^B\to\{g\mid g:\lambda\to\kappa\}$ by $T(f)\defeq u\circ f\circ v^{-1}$. Its inverse sends $g:\lambda\to\kappa$ to $u^{-1}\circ g\circ v$, so $T$ is bijective by Theorem \ref{thm:bijective-iff-invertible-function}.\\
Hence $\kappa^\lambda=\abs{A^B}$ by Proposition \ref{prop:cardinality-canonical-and-bijection-invariant}(ii).
\end{enumerate}
\hlt{Structural Commentary}: Chosen bijections transport comparison and all three cardinal operations.\\
\hlt{Key Insight}: Bijections act coordinatewise on tags, pairs, and functions.
\end{proof}
\end{proposition}

\begin{definition} \label{def:continuum-cardinal} \hlt{Continuum Cardinal}\\
The \hlt{continuum cardinal} is $\mathfrak{c}\defeq2^{\aleph_0}$.
\end{definition}

\begin{proposition} \label{prop:continuum-cardinal-models} \hlt{Models of the Continuum Cardinal}\\
$\mathfrak{c}=\abs{\{0,1\}^{\N}}=\abs{\mathcal{P}(\N)}$.
\begin{proof}
The map $q:J_2\to\{0,1\}$ defined by $q(k)\defeq k-1$ is bijective. Thus $\abs{\{0,1\}}=2$ and $\abs{\N}=\aleph_0$ by Proposition \ref{prop:finite-cardinality-and-aleph-zero-agreement}(i),(iii).\\
Hence $\mathfrak{c}=\abs{\{0,1\}^{\N}}$ by Proposition \ref{prop:cardinal-operations-on-representatives}(v). The relation $\{0,1\}^{\N}\sim_{\#}\mathcal{P}(\N)$ holds by Proposition \ref{prop:binary-sequences-model-power-set}.\\
Hence $\abs{\{0,1\}^{\N}}=\abs{\mathcal{P}(\N)}$ by Proposition \ref{prop:cardinality-canonical-and-bijection-invariant}(ii).\\
\hlt{Structural Commentary}: Finite and countable representatives reduce cardinal exponentiation to binary functions and subsets.\\
\hlt{Key Insight}: A binary function is the characteristic function of one subset.
\end{proof}
\end{proposition}

\begin{corollary} \label{cor:continuum-strictly-exceeds-aleph-zero} \hlt{The Continuum Strictly Exceeds $\aleph_0$}\\
$\aleph_0<\mathfrak{c}$.
\begin{proof}
The map $j:\N\to\mathcal{P}(\N)$, $j(n)\defeq\{n\}$, is injective since $j(m)=j(n)\Rightarrow m=n$.\\
No bijection $\N\to\mathcal{P}(\N)$ exists by Theorem \ref{thm:cantor-power-set-theorem}. Therefore $\abs{\N}<\abs{\mathcal{P}(\N)}$ by Proposition \ref{prop:cardinal-operations-on-representatives}(ii).\\
The identities $\abs{\N}=\aleph_0$ and $\abs{\mathcal{P}(\N)}=\mathfrak{c}$ hold by Proposition \ref{prop:finite-cardinality-and-aleph-zero-agreement}(iii) and Proposition \ref{prop:continuum-cardinal-models}, respectively. Hence $\aleph_0<\mathfrak{c}$.\\
\hlt{Structural Commentary}: Singleton embedding gives weak comparison, while Cantor's theorem rules out equality.\\
\hlt{Key Insight}: The diagonal obstruction makes the continuum strictly larger.
\end{proof}
\end{corollary}
